\section{Design of \tool{}}
\label{sec:design}

\subsection{Principles}

\tool{} is a web application and Python library (FastAPI back end, React front end, MIT licence). Three principles
shape it:
\begin{enumerate}
\item \textbf{Protocol first.} A question is turned into a review protocol before anything is searched: the review
  type, the parts of the question (concepts) with their synonyms, the populations or places every study must name,
  and the eligibility criteria. The tool proposes a protocol from the question's own words; in the guided review
  the user confirms or edits each part. The protocol is stored with the results.
\item \textbf{A bounded role for the model.} A model may propose synonyms and search aspects, screen titles and
  abstracts against the protocol's criteria, extract study details and write the synthesis. It never writes search
  strings, and the rule-based steps that decide inclusion (duplicate merging, eligibility and ranking) do not read
  its output: they match studies against the concepts parsed from the question by fixed rules, so the same records
  give the same candidates with any model or none.
\item \textbf{Everything stored, selection replayable.} Each run is stored as a \emph{corpus}: every record the
  searches returned (before and after merging), the search plan and strings, the protocol, every eligibility and
  screening decision with its reason and the model's name, the selection protocol version and the code commit,
  under a content hash. A replay function re-runs the selection from the stored records alone, with no network
  and no model, and checks that it reproduces the included studies, in order, with the same hash.
\end{enumerate}
Figure~\ref{fig:pipeline} shows the pipeline. The full selection rules are published in the repository
(\texttt{docs/RANKING.md}, generated from the code) and summarised below; they are those of selection protocol
\texttt{citelume-selection/2.5}, the version evaluated.

\begin{figure}[t]
\centering
\begin{tikzpicture}[
  box/.style={draw, rounded corners=2pt, align=center, font=\scriptsize, minimum height=12mm, text width=25mm},
  model/.style={box, fill=orange!15},
  rule/.style={box, fill=blue!8},
  store/.style={draw, rounded corners=2pt, align=center, font=\scriptsize, fill=gray!15, text width=112mm,
    minimum height=8mm},
  arr/.style={-{Stealth[length=2mm]}, thick}]
\node[rule]  (q)       at (0,0)    {Question $\rightarrow$ review protocol (user confirms)};
\node[model] (plan)    at (3.6,0)  {Search plan: concepts, synonyms, aspects};
\node[rule]  (strings) at (7.2,0)  {Search strings compiled by rules (+ MeSH)};
\node[rule]  (search)  at (10.8,0) {Database searches; citation searching};
\node[rule]  (merge)   at (10.8,-2.2) {Merge duplicates; eligibility rules; 100-point score};
\node[model] (screen)  at (7.2,-2.2)  {Model screens titles and abstracts against criteria};
\node[rule]  (incl)    at (3.6,-2.2)  {Included studies: every eligible, up to the limit};
\node[model] (synth)   at (0,-2.2)    {Extraction, synthesis and checks};
\node[store] (corpus)  at (5.4,-4.1)  {Stored corpus: records, plan, strings, every decision and reason, model names,
  versions, content hash $\rightarrow$ \emph{the selection replays without search or model}};
\draw[arr] (q) -- (plan); \draw[arr] (plan) -- (strings); \draw[arr] (strings) -- (search);
\draw[arr] (search) -- (merge); \draw[arr] (merge) -- (screen); \draw[arr] (screen) -- (incl);
\draw[arr] (incl) -- (synth);
\draw[arr, dashed, gray] (merge.south) -- (merge.south |- corpus.north);
\draw[arr, dashed, gray] (screen.south) -- (screen.south |- corpus.north);
\draw[arr, dashed, gray] (incl.south) -- (incl.south |- corpus.north);
\end{tikzpicture}
\caption{The \tool{} pipeline. Orange steps may use a model; blue steps follow fixed rules. The rule-based steps that
decide which studies are candidates do not read the model's plan. Everything is stored in a corpus from which the
selection can be replayed.}
\label{fig:pipeline}
\end{figure}

\subsection{Search}

The question is reduced to one to three concepts that every relevant record must be about, each with synonyms;
a population or place named in the question is also required. With a model configured, the model proposes the
concepts, synonyms and four to six search aspects (for example mechanisms, adverse effects), and its concepts are
accepted only if the leading one is anchored in the question's words; otherwise fixed rules built for questions
written as review titles derive them. Search strings are then compiled by rules, never written by the model. For
each concept, PubMed is asked once how it maps the words to its controlled vocabulary, and the returned MeSH
headings and pharmacological actions are added as alternatives \citep{lefebvre2023searching}. Each database (PubMed
with its Europe PMC mirror, Semantic Scholar with OpenAlex as a fallback, bioRxiv/medRxiv and arXiv) receives a core
search (the concepts), a reviews search and one search per aspect. For systematic-style reviews the core PubMed
search can be paged deeper (500 records in the evaluation). After a first selection, the reference lists and citing
papers of up to 60 selected studies are fetched from Semantic Scholar (records linked to at least two are added),
and the reference lists of on-topic reviews set aside by the eligibility rules are added, as records identified by
other methods \citep{page2021prisma}. The strings, the number of records each PubMed search matched and how many
were retrieved are stored and shown.

\subsection{Selection}

Records for the same paper are merged by DOI, PubMed ID, arXiv ID or normalised title. Eligibility rules then
remove, each under its own reason: records outside a population or place the question names; records whose main
concept appears nowhere in the title, abstract or (for PubMed records) MeSH indexing; records that are not research
reports (corrections, retraction notices, protocols without results, meeting abstracts); and records outside the
protocol's years or designs. In a systematic-style review, reviews of all kinds are set aside and listed for the
discussion. Each remaining record receives a published 100-point score: relevance to the question (50 points:
concept matches in title and abstract, and BM25 wording overlap ranked among the candidates), study type (20 points,
from the database's publication type where available, otherwise from wording), citations per year ranked among the
candidates (15), recency (10) and the number of independent searches that found it (5).

For a quick answer the 20 highest-scoring studies are chosen with simple coverage rules (each search aspect is
represented; reviews limited to a share of places). For a systematic-style review there are no quotas: every
eligible primary study whose title or abstract addresses every part of the question is a candidate, and records
that match the main concepts but not every part go to model screening rather than being excluded by word matching.
The model then screens candidates in score order, in batches of 12, answering each protocol criterion as yes, no
or unclear for each record from its title and the first 420 characters of its abstract; a record is excluded only
if a criterion is clearly not met (``unclear'' never excludes). Screening continues until the study limit is filled
(150 in the evaluation), up to the larger of 240 records or twice the limit per review. Each decision is stored
with its reason and can be overturned by the user, with a recorded reason. Every exclusion is counted at its PRISMA
2020 stage for the flow diagram \citep{page2021prisma}.

\subsection{After selection}

A model extracts each included study's design, population, results and stated limitations (from open-access full
text where available) and writes a synthesis in which every sentence cites the studies it rests on; rule-based
checks compare figures and citations with the extracted details. Risk of bias is suggested only for studies with
open full text, each judgement resting on a quotation checked against the text and, for randomised trials, on the
RoB 2 algorithm \citep{sterne2019rob2}; it is presented as a suggestion for the reviewer and was not evaluated
against published judgements (Section~\ref{sec:limits}). This paper evaluates search and selection, the
repeatability of each step, and the accuracy of the written claims.

\section{Evaluation}
\label{sec:eval}

The evaluation protocol was written and committed on 18 November 2026, before any benchmark review was loaded or
run; changes were made only by dated amendments (Supplementary Material S2). The search methods are reported
following PRISMA-S \citep{rethlefsen2021prismas} (Supplementary Material S1). All code, stored corpora and results
are available (Section~\ref{sec:availability}).

\subsection{Benchmark}

We used the 40 Cochrane intervention reviews of CLEF eHealth TAR 2019, Task~2 \citep{kanoulas2019clef}, each with
the PubMed IDs of the studies the review included after full-text screening. Intervention reviews were chosen
because they ask effect questions, the kind \tool{}'s systematic-style protocol is built for; diagnostic and
prognostic topics were not used. Two reviews listed in both parts of the track were counted once, and five were
excluded by rules fixed beforehand (fewer than two included studies resolving to a PubMed record), leaving 33
reviews with 785 included studies (141 in the development set and 644 in the test set). The reviews were sorted by ID, shuffled with a fixed seed (20261118), and the first third formed the
development set (11 reviews); the other 22 formed the test set. Development reviews were used to change the
pipeline; test reviews were run only after the protocol was fixed.

\subsection{Runs}

Each review's title, unchanged, was the question; searches were limited to records published up to the year of the
review version CLEF used. Each test review was run in the systematic-style mode (Deep Research, PubMed core search
500 records deep, citation searching and review reference lists on, study limit 150) and in the quick-answer mode.
The model for planning, screening and analysis was NVIDIA Nemotron 3 Ultra (\url{nvidia/nemotron-3-ultra-550b-a55b})
through OpenRouter, on its free endpoint for the systematic-style runs and its paid endpoint for the quick-answer
runs (total cost 0.77~USD). All 22 searches in each mode were planned by the model. Every run went through a batch
runner that stores the corpus, result, timings, model and code commit, and fails a run in which more than 5\% of the
records due for screening were not screened.

\subsection{Freezing the measured code}

Search and selection were frozen at selection protocol 2.5 (code commit 47e40ef, git tag
\texttt{eval-freeze-2.5}) on 23 November 2026, after the systematic-style test runs and before any further run
(protocol amendment, Supplementary Material S2). The hashes of the frozen files and functions are recorded and
checked by an automated test, so any later change to search or selection fails the test suite until it is
recorded as a post-freeze change. One post-freeze fix was made (a bioRxiv connector error that did not occur in any
evaluation run).

\subsection{Measures}

The primary measure was \emph{pool recall}: the share of a review's included studies present in the candidate pool
after duplicate merging. Secondary measures were the share in the top 20 and top 60 of the screening order, the
share selected after model screening, screening sensitivity (included studies kept by the model, of those it
screened) and precision, and the loss of included studies at each stage (search, merging, eligibility, screening and
study limit). A study counted as found when a record carried its PubMed ID or DOI (case-insensitive) or, for records
with neither, the same normalised title. Each measure was computed per review and summarised as the mean with a 95\%
bootstrap confidence interval over reviews (10,000 resamples, seed 20261118) and the median, over all reviews and
over the reviews with at most 60 included studies. Baselines ranked the same pools by BM25 against the question, by
citation count alone, and at random.

\subsection{Repeatability and sensitivity}

On the stored test corpora we measured (E2a) whether the selection replays exactly from the stored records;
(E2b) the agreement of model screening with itself, by screening a sample of 1,638 records (up to 96 per review,
every screened included study first) three more times with the same model and prompt; (E2c) the variation in
records found when each question is searched twice more; (E2d) the effect of model planning, by repeating the
searches with rule-based planning and re-planning each question three times with the model; and (E3) the
sensitivity of the selection to its constants, by setting each score weight and cap to half and one and a half
times its value and comparing the selections (Jaccard overlap) and recall.

\subsection{Component checks}

We checked the study-type classifier against the publication types PubMed indexers assigned to 26,114 records from
the evaluation runs, with those types removed so the classifier had to rely on wording, as it does for most
preprint and some Semantic Scholar records. We sampled 100 duplicate merges from the test runs (1,351 merges on title
alone and 99 merges of records sharing an identifier with differing titles), weighted to the harder kinds, for a
person to judge whether each pair was the same study \PENDING{merge check result}. To measure the accuracy of the
written synthesis, 200 cited sentences were drawn at random (seed 20261118) from the 881 shown to readers in 33
model-written analyses (the 22 test questions and the 11 development questions, quick-answer mode), and a person
judged each against the abstracts of the studies it cites as supported, partly supported, not supported or cannot
tell, using a written guide and blinded to how each analysis was produced; a second person judged 50 of them
independently \PENDING{annotators and agreement}.

\subsection{Deviations from the protocol}

(1) The rule for which review version to use was amended before any run, because CLEF used nine reviews whose only
PubMed version is dated March 2019. (2) The ranking change in selection protocol 2.5 was measured on the
development set without the model before the test runs. (3) One review was re-run after dropped connections lost
two screening batches; in the final runs two reviews each have 12 records the model did not screen (under the 5\%
threshold). (4) Development runs used the paid endpoint for 9 of 11 reviews; test systematic-style runs used the
free endpoint and quick-answer runs the paid one. (5) Search and selection were frozen after the systematic-style
test results had been seen; no change was made to them on account of those results.
